\documentclass[11pt,a4paper]{article}
\usepackage[T1]{fontenc}
\usepackage[utf8]{inputenc}
\usepackage{lmodern}
\usepackage[margin=25mm]{geometry}
\usepackage{microtype}
\usepackage[hidelinks]{hyperref}
\setlength{\parindent}{0pt}
\setlength{\parskip}{0.75em}
\pagestyle{empty}
\begin{document}
\textbf{To the Editors of Statistics: A Journal of Theoretical and Applied Statistics}

Dear Editors,

Please consider my manuscript, ``Finite-sample confidence sets for nonsmooth
convex-to-concave transitions,'' as an original article in \emph{Statistics}.

The paper addresses an uncertainty problem in shape-constrained regression:
which transition locations can be ruled out when a curve changes from convex
to concave, but the change may involve a kink, a jump, or a flat region?
A point estimate does not answer this question, and inference built for a
smooth, unique inflection need not remain valid in these settings.

The proposed method, shape-contrast inversion (SCI), combines simultaneous
bounds for local chord contrasts with a direct exclusion rule. Its target is
the full set of transitions admitted by shape-constrained continuations of
the fixed-design mean vector. Under the stated Gaussian assumptions, the
confidence set covers this target uniformly at finite sample sizes.
Extensions address unknown constant noise, bounded variance heterogeneity,
and independent replicate curves with within-curve dependence.

The contribution is not a new point estimator or the first method for
inflection-point inference. It is a directly computable confidence procedure
for a set-valued transition target that permits nonsmoothness and nonuniqueness.
The paper also gives localization bounds under explicit contrast-margin
conditions, without claiming minimax optimality.

An 80,000-response simulation audit checks full-set coverage. Comparisons
include a published smooth-bootstrap method and a conservative baseline with
a matched coverage guarantee. They report both useful localization on
nonsmooth signals and weak-signal cases where SCI remains wide. LIDAR and
replicated DNase examples illustrate the dependence on noise assumptions.
Code and numerical artifacts are archived in Zenodo at
\url{https://doi.org/10.5281/zenodo.22338783} (release sci-v1.0.0).
A matrix-free implementation supports the statistical contribution.

I believe this combination of finite-sample theory, explicit limitations, and
applied examples fits the journal's theoretical and applied scope. I am the
sole author and conducted this research independently. This manuscript is not
under consideration by another journal. The work received no funding, and I report no competing
interests. The manuscript includes a declaration of generative AI use.

Thank you for considering the manuscript.

Sincerely,\\
Saveliy Baturin\\
Independent Researcher\\
\href{mailto:saveliy.xo@gmail.com}{saveliy.xo@gmail.com}
\end{document}
